\documentclass[aspectratio=169]{beamer}

\usetheme{Madrid}
\usepackage{booktabs}
\usepackage{xcolor}
\setbeamertemplate{navigation symbols}{}

\title[Rules, provenance, and raw-Sentinel water]{Interpretable rule splits, a STACD provenance proposal, per-fortnight water, and safe defaults}
\author[Salil Gujar]{Salil Gujar \\ \footnotesize M.Tech CSE \quad Entry No.\ 2025MCS2106}
\date{July 2026}

\begin{document}

\frame{\titlepage}

% ------------------------------------------------------------------
\begin{frame}{Where we are, and what week 9 tackled}
\begin{columns}[T]
\begin{column}{0.5\textwidth}
\emph{Done so far}
\begin{itemize}\small
  \item A 4-class base map over India at 10\,m with a living hierarchy: split or add a class, mark a
  few example polygons, and retrain that node on the fly.
  \item A git-backed zoo of model and dataset cards, with save and resume, merge, base and year
  selection, a GeoTIFF export, and site testing on acacia, tea, and mining.
\end{itemize}
\end{column}
\begin{column}{0.5\textwidth}
\emph{This week (a chosen subset)}
\begin{itemize}\small
  \item Split a class by an interpretable rule on indices, not only by a trained model.
  \item Read the STACD paper, and propose how a provenance record for our outputs would look.
  \item Answer water seasonality on raw Sentinel, one fortnight at a time.
  \item Add safe defaults: a bounding-box size cap and a training-time benchmark.
  \item Link the valid model families to the inference data the user picks.
\end{itemize}
\end{column}
\end{columns}
\end{frame}

% ------------------------------------------------------------------
\begin{frame}{Splitting a class by a rule, not a model}
\begin{itemize}
  \item Until now, dividing a class meant marking data and training a classifier. Sometimes the
  boundary is simply a threshold on a known index, for example dense against sparse vegetation by
  annual NDVI.
  \item A rule split lets the user write that directly: if the expression holds the pixel takes one
  class, otherwise it takes the default. No data to mark, no training.
  \item The indices come from a fixed registry, so the expressions stay interpretable and safe. We
  check every expression against a whitelist of variables and operators before it reaches Earth
  Engine.
  \item Because the indices are computed in Earth Engine, a rule split renders as the same crisp
  10\,m tile map as a linear model. It is stored on the node, travels with the saved project, and
  gets a card in the zoo like any other model.
\end{itemize}
\end{frame}

% ------------------------------------------------------------------
\begin{frame}{The rule variable registry}
The indices a rule can use, all computed live from Sentinel and SRTM over the chosen area and year.

\vspace{3mm}
\begin{center}
\begin{tabular}{ll}
\toprule
Group & Variables \\
\midrule
Vegetation & ndvi\_annual, ndvi\_kharif, ndvi\_rabi \\
Water & ndwi, mndwi \\
Built and soil & ndbi, bsi \\
Terrain (SRTM) & slope, elevation \\
\bottomrule
\end{tabular}
\end{center}
\end{frame}

% ------------------------------------------------------------------
\begin{frame}{STACD: what I understood from the paper}
This is my reading of the STACD paper, presented so we can confirm it together.
\begin{itemize}\small
  \item STAC describes an individual geospatial asset well, but its lineage is shallow: a
  \texttt{derived\_from} field that names a parent without saying which algorithm, version, or
  parameters produced it.
  \item STACD extends STAC with a directed acyclic graph and five classes: a DAG, a Dataset Type
  and an Algorithm Type as the node kinds, an Algorithm Instance that records a specific version and
  its code, and a Dataset Instance that extends a STAC Item with the producing algorithm and its
  inputs.
  \item The payoff is provenance and selective recomputation: if a model or an input dataset
  changes, you can see exactly which downstream outputs need to be rebuilt.
  \item The paper gives a reference implementation on Apache Airflow.
\end{itemize}
\end{frame}

% ------------------------------------------------------------------
\begin{frame}{How our output could map onto STACD}
Something like this could work for us. A classified output would carry two things.
\begin{itemize}\small
  \item A stack specification: a STAC Item for the LULC raster, with the bounding box, the geometry,
  the class legend with colours, and asset links to the GeoTIFF and tile endpoints.
  \item A STACD specification: the DAG plus the dataset and algorithm nodes, and one algorithm
  instance for every live model, rule split, and merge, each pointing at where its artifact lives in
  the zoo.
  \item The input set: the class hierarchy, the ordered operations that built it, and the artifact
  pointers, carried as the inputs of the producing algorithm. This is the JSON we already keep, which
  would sit inside the STACD record as the input set that produced the output.
  \item The appeal is that it reuses the cards we already have, so it needs no new source of truth,
  and it stays a cheap metadata step with no extra Earth Engine work.
\end{itemize}
\end{frame}

% ------------------------------------------------------------------
\begin{frame}{Water seasonality, one fortnight at a time}
\begin{itemize}\small
  \item The annual embedding is a whole-year signal, so it cannot say which pixels held water on a
  given fortnight. For that we go back to raw Sentinel: Sentinel-1 radar and Sentinel-2 optical,
  composited around a target date, with the standard water indices.
  \item We train a linear water against non-water model offline on the shared seasonal-water
  polygons, sampled at each polygon's own date, then replay it as Earth-Engine band math for any
  chosen fortnight. This keeps it interactive, the same trick the base map uses, and avoids the
  memory limit that pushed the earlier flood pipeline to batch export.
  \item Early results are good enough to work with for now, and we can tighten them as we add more
  polygons and dates.
\end{itemize}
\end{frame}

% ------------------------------------------------------------------
\begin{frame}{A size cap on the drawn area}
\begin{itemize}
  \item A user can draw any rectangle, and a large one can quietly turn into a very slow request or
  a large download. We now check the area before running.
  \item The three paths cost differently, so the caps differ: the Earth-Engine tile map is generous
  because tiles render on demand, the GeoTIFF export is smaller because that download is size
  limited, and the Tessera path is capped by tile count because each tile is a large local
  download.
  \item The caps are plain environment settings, so a server administrator can tune them for the
  hardware. The interface reports the area and refuses to run past the hard limit, with a readable
  reason.
\end{itemize}
\end{frame}

% ------------------------------------------------------------------
\begin{frame}{How long does training take}
A benchmark separates the two costs so an administrator can size the server. Sampling embeddings
over Earth Engine is network bound and usually dominant; the linear fit is cheap. Random Forest is a
local reference only.
\vspace{3mm}
\centering
\begin{tabular}{lrrr}
\toprule
Area (km$^2$) & Points at 10\% & Sample (s) & Total, linear (s) \\
\midrule
1 & 1{,}000 & 8.1 & 8.1 \\
4 & 4{,}000 & 16.1 & 16.2 \\
25 & 25{,}000 & 72.7 & 72.9 \\
100 & 100{,}000 & 274.5 & 275.7 \\
\bottomrule
\end{tabular}
\vspace{3mm}
\begin{flushleft}\small
Measured on a Windows laptop. The numbers feed a saved profile and a live estimate endpoint, and
they are the reference for setting the area caps above.
\end{flushleft}
\end{frame}

% ------------------------------------------------------------------
\begin{frame}{Linking models to the inference data}
\begin{itemize}
  \item When a user trains, the valid model families now follow the data they picked. Earth Engine
  can only run linear models, because only a linear model replays as band math, so the Alpha Earth
  list is linear only.
  \item A Tessera run is a local computation, so it is not limited to linear. It adds Random Forest,
  and XGBoost when installed, and reserves a slot for an object-detection family we have not built
  yet.
  \item The training panel reads this list, so the user is only offered models that fit their data,
  and a non-linear model on Alpha Earth is refused with a clear reason.
  \item This is the linkage advised earlier: the dataset a user works with determines the models
  available for it, and the set is easy to extend.
\end{itemize}
\end{frame}

% ------------------------------------------------------------------
\begin{frame}{Improving the acacia split}
Acacia against non-acacia is our hardest split, a species distinction between similar trees. The
options we have actually measured on it:
\begin{itemize}
  \item More labelled crowns across sites and years. Data spread is the proven lever: pooling several
  years moved this split from 0.635 to 0.745 on years the model had never seen.
  \item Tessera as the feature source, tested at about 0.73 on the same split. With the linkage
  above it can now also carry a non-linear model such as Random Forest, which the 64-dimensional
  Alpha Earth path cannot.
  \item The automatic bake-off, already run on this split, where Ridge came out ahead at 0.731.
\end{itemize}
\end{frame}

% ------------------------------------------------------------------
\begin{frame}{Where the product stands}
\begin{itemize}
  \item The full loop works: paint a map, grow a scheme with splits, rule splits, adds and merges,
  retrain, browse and publish the zoo, save and resume a project, and export a GeoTIFF. This is a
  demo-ready product.
  \item What stands between a demo and a hosted service is mostly packaging: a service-account key
  for headless Earth Engine, a container, pinned dependencies, and a basic access limit so the
  compute is not open to anyone. The code is modular, so these are additions, not rewrites.
  \item The water model is new and is only lightly wired into the interface so far.
\end{itemize}
\end{frame}

% ------------------------------------------------------------------
\begin{frame}
\centering
\vspace{1.5cm}
{\Large Thank you}
\vspace{0.5cm}


\end{frame}

\end{document}
